\documentclass[10pt,a4paper]{amsart}
\usepackage{amsmath,amssymb,amsthm,graphicx}
% Retained preamble for the Stacks Project chapter snapshot in this directory.
% The chapter's own preamble.tex is a copy of the shared Stacks Project preamble; the block
% below is sliced MECHANICALLY, line for line, from it. Nothing was added, removed or
% redefined; the wrapper class, the xr-hyper/hyperref wiring and the externaldocument list
% are omitted because the assembler supplies the document class.
%
% lines taken from preamble.tex: 3, 5, 9, 11, 12, 156, 157, 158, ... (51 lines total)
%
% package coverage restored from the chapter's own preamble (2026-10-13): xy, xy-2cell

\usepackage[all]{xy}
\xyoption{2cell}
\UseAllTwocells

\newenvironment{reference}{\comment}{\endcomment}
\newenvironment{slogan}{\comment}{\endcomment}
\newenvironment{history}{\comment}{\endcomment}
\theoremstyle{plain}
\newtheorem{theorem}[subsection]{Theorem}
\newtheorem{proposition}[subsection]{Proposition}
\newtheorem{lemma}[subsection]{Lemma}
\theoremstyle{definition}
\newtheorem{definition}[subsection]{Definition}
\newtheorem{example}[subsection]{Example}
\newtheorem{exercise}[subsection]{Exercise}
\newtheorem{situation}[subsection]{Situation}
\theoremstyle{remark}
\newtheorem{remark}[subsection]{Remark}
\newtheorem{remarks}[subsection]{Remarks}
\numberwithin{equation}{subsection}
\def\lim{\mathop{\mathrm{lim}}\nolimits}
\def\colim{\mathop{\mathrm{colim}}\nolimits}
\def\Spec{\mathop{\mathrm{Spec}}}
\def\Hom{\mathop{\mathrm{Hom}}\nolimits}
\def\Ext{\mathop{\mathrm{Ext}}\nolimits}
\def\SheafHom{\mathop{\mathcal{H}\!\mathit{om}}\nolimits}
\def\SheafExt{\mathop{\mathcal{E}\!\mathit{xt}}\nolimits}
\def\Sch{\mathit{Sch}}
\def\Mor{\mathop{\mathrm{Mor}}\nolimits}
\def\Ob{\mathop{\mathrm{Ob}}\nolimits}
\def\Sh{\mathop{\mathit{Sh}}\nolimits}
\def\NL{\mathop{N\!L}\nolimits}
\def\CH{\mathop{\mathrm{CH}}\nolimits}
\def\proetale{{pro\text{-}\acute{e}tale}}
\def\etale{{\acute{e}tale}}
\def\QCoh{\mathit{QCoh}}
\def\Ker{\mathop{\mathrm{Ker}}}
\def\Im{\mathop{\mathrm{Im}}}
\def\Coker{\mathop{\mathrm{Coker}}}
\def\Coim{\mathop{\mathrm{Coim}}}
\DeclareMathSymbol{\boxtimes}{\mathbin}{AMSa}{"02}
\def\QCohstack{\mathcal{QC}\!\mathit{oh}}
\def\Cohstack{\mathcal{C}\!\mathit{oh}}
\def\Spacesstack{\mathcal{S}\!\mathit{paces}}
\def\Quotfunctor{\mathrm{Quot}}
\def\Hilbfunctor{\mathrm{Hilb}}
\def\Curvesstack{\mathcal{C}\!\mathit{urves}}
\def\Polarizedstack{\mathcal{P}\!\mathit{olarized}}
\def\Complexesstack{\mathcal{C}\!\mathit{omplexes}}
\def\Pic{\mathop{\mathrm{Pic}}\nolimits}
\def\Picardstack{\mathcal{P}\!\mathit{ic}}
\def\Picardfunctor{\mathrm{Pic}}
\def\Deformationcategory{\mathcal{D}\!\mathit{ef}}
\usepackage{cleveref}
\crefname{originallemma}{original lemma}{original lemmas}
\Crefname{originallemma}{Original lemma}{Original lemmas}
\makeatletter
\newlabel{equation-system-coequalizers}{{0AGJ}{source}{}{sourcepointer-1}{}}
\newlabel{equation-system-coequalizers@cref}{{[originallemma][0][]0AGJ}{[1][0][]source}{}{}{}}
\makeatother
\crefname{originallemma}{original lemma}{original lemmas}
\Crefname{originallemma}{Original lemma}{Original lemmas}
\makeatletter
\newlabel{situation-coequalizer-glue}{{0AGH}{source}{}{sourcepointer-2}{}}
\newlabel{situation-coequalizer-glue@cref}{{[originallemma][0][]0AGH}{[1][0][]source}{}{}{}}
\makeatother
\crefname{originallemma}{original lemma}{original lemmas}
\Crefname{originallemma}{Original lemma}{Original lemmas}
\makeatletter
\newlabel{spaces-cohomology-lemma-Artin-Rees}{{07UL}{source}{}{sourcepointer-3}{}}
\newlabel{spaces-cohomology-lemma-Artin-Rees@cref}{{[originallemma][0][]07UL}{[1][0][]source}{}{}{}}
\makeatother
\crefname{originallemma}{original lemma}{original lemmas}
\Crefname{originallemma}{Original lemma}{Original lemmas}
\makeatletter
\newlabel{spaces-cohomology-lemma-ML-cohomology-powers-ideal}{{08AY}{source}{}{sourcepointer-4}{}}
\newlabel{spaces-cohomology-lemma-ML-cohomology-powers-ideal@cref}{{[originallemma][0][]08AY}{[1][0][]source}{}{}{}}
\makeatother
\crefname{originallemma}{original lemma}{original lemmas}
\Crefname{originallemma}{Original lemma}{Original lemmas}
\makeatletter
\newlabel{spaces-cohomology-lemma-acc-coherent}{{07UJ}{source}{}{sourcepointer-5}{}}
\newlabel{spaces-cohomology-lemma-acc-coherent@cref}{{[originallemma][0][]07UJ}{[1][0][]source}{}{}{}}
\makeatother
\crefname{originallemma}{original lemma}{original lemmas}
\Crefname{originallemma}{Original lemma}{Original lemmas}
\makeatletter
\newlabel{spaces-cohomology-lemma-power-ideal-kills-sheaf}{{07UK}{source}{}{sourcepointer-6}{}}
\newlabel{spaces-cohomology-lemma-power-ideal-kills-sheaf@cref}{{[originallemma][0][]07UK}{[1][0][]source}{}{}{}}
\makeatother
\crefname{originallemma}{original lemma}{original lemmas}
\Crefname{originallemma}{Original lemma}{Original lemmas}
\makeatletter
\newlabel{spaces-cohomology-section-theorem-formal-functions}{{08AU}{source}{}{sourcepointer-7}{}}
\newlabel{spaces-cohomology-section-theorem-formal-functions@cref}{{[originallemma][0][]08AU}{[1][0][]source}{}{}{}}
\makeatother
\crefname{originallemma}{original lemma}{original lemmas}
\Crefname{originallemma}{Original lemma}{Original lemmas}
\makeatletter
\newlabel{spaces-properties-lemma-quasi-compact-affine-cover}{{03H6}{source}{}{sourcepointer-8}{}}
\newlabel{spaces-properties-lemma-quasi-compact-affine-cover@cref}{{[originallemma][0][]03H6}{[1][0][]source}{}{}{}}
\makeatother
\title{The coequalizer system X\_n is essentially constant with limit X for quasi-compact X}
\author{The Stacks Project Authors}
\date{stacks-project:spaces-pushouts.tex:local-snapshot-20261012; GNU FDL 1.2 or later}
\begin{document}
\maketitle
\noindent\textit{Editorial note.} Statement span is lines 3066-3072 of sources/stacks-spaces-pushouts/pinned-original-source.tex; proof span is lines 3074-3185, the complete author proof (110 body lines).\par\medskip
\noindent\textit{Editorial note.} Navigation stubs: 8 cited-but-not-included labels are supplied as non-load-bearing pointers with their pinned tags where the snapshot tag index has them.\par\medskip
\noindent\textit{Editorial note.} extra\_packages is empty on purpose: the retained preamble already loads xy with options and a second declaration clashes in the assembler wrapper.\par\medskip
\noindent\textit{Editorial note.} Text is retained verbatim, including the author's own typos and placeholders.\par\medskip
\noindent\textit{Editorial note.} Queue index 2 of queue\_spaces-pushouts.json (pb=4576, stmt 3066-3072, proof 3074-3185).\par\medskip
\noindent\textit{Editorial note.} spaces-pushouts wave-1: fresh adjudication accept (2026-10-14); six-check dispatched pool (§2 row appended 02:4x).\par\medskip
\noindent\textit{Allowed standard prerequisites (editorial).}\begin{itemize}
\item Explicit prerequisite boundary. The cited inputs are prerequisites; the argument in the retained span is not replaced by any cited result.
\end{itemize}\medskip
\ifdefined\hypertarget\hypertarget{sourcepointer-1}{}\fi
\noindent\textit{Source reference.} Original lemma 0AGJ: sources/stacks-spaces-pushouts/pinned-original-source.tex, source lines 3060--3062. Cited by the retained statement or proof but not part of the collected spans; supplied only as a navigation pointer. The argument does not depend on its content.\par\medskip
\ifdefined\hypertarget\hypertarget{sourcepointer-2}{}\fi
\noindent\textit{Source reference.} Original lemma 0AGH: sources/stacks-spaces-pushouts/pinned-original-source.tex, source lines 3021--3026. Cited by the retained statement or proof but not part of the collected spans; supplied only as a navigation pointer. The argument does not depend on its content.\par\medskip
\ifdefined\hypertarget\hypertarget{sourcepointer-3}{}\fi
\noindent\textit{Source reference.} Original lemma 07UL: sources/stacks-spaces-pushouts/pinned-original-source.tex, source lines 3185--3185. Cited by the retained statement or proof but not part of the collected spans; supplied only as a navigation pointer. The argument does not depend on its content.\par\medskip
\ifdefined\hypertarget\hypertarget{sourcepointer-4}{}\fi
\noindent\textit{Source reference.} Original lemma 08AY: sources/stacks-spaces-pushouts/pinned-original-source.tex, source lines 3185--3185. Cited by the retained statement or proof but not part of the collected spans; supplied only as a navigation pointer. The argument does not depend on its content.\par\medskip
\ifdefined\hypertarget\hypertarget{sourcepointer-5}{}\fi
\noindent\textit{Source reference.} Original lemma 07UJ: sources/stacks-spaces-pushouts/pinned-original-source.tex, source lines 3185--3185. Cited by the retained statement or proof but not part of the collected spans; supplied only as a navigation pointer. The argument does not depend on its content.\par\medskip
\ifdefined\hypertarget\hypertarget{sourcepointer-6}{}\fi
\noindent\textit{Source reference.} Original lemma 07UK: sources/stacks-spaces-pushouts/pinned-original-source.tex, source lines 3185--3185. Cited by the retained statement or proof but not part of the collected spans; supplied only as a navigation pointer. The argument does not depend on its content.\par\medskip
\ifdefined\hypertarget\hypertarget{sourcepointer-7}{}\fi
\noindent\textit{Source reference.} Original lemma 08AU: sources/stacks-spaces-pushouts/pinned-original-source.tex, source lines 3185--3185. Cited by the retained statement or proof but not part of the collected spans; supplied only as a navigation pointer. The argument does not depend on its content.\par\medskip
\ifdefined\hypertarget\hypertarget{sourcepointer-8}{}\fi
\noindent\textit{Source reference.} Original lemma 03H6: sources/stacks-spaces-pushouts/pinned-original-source.tex, source lines 3185--3185. Cited by the retained statement or proof but not part of the collected spans; supplied only as a navigation pointer. The argument does not depend on its content.\par\medskip
\begin{lemma}
\label{lemma-essentially-constant}
In Situation \ref{situation-coequalizer-glue} assume $X$ quasi-compact.
In (\ref{equation-system-coequalizers}) for all $n$ large enough, there
exists an $m$ such that $X_n \to X_{n + m}$ factors through a
closed immersion $X \to X_{n + m}$.
\end{lemma}

\begin{proof}
Let's look a bit more closely at the construction of $X_n$
and how it changes as we increase $n$. We have
$X_n = \underline{\Spec}(\mathcal{A}_n)$
where $\mathcal{A}_n$ is the equalizer of $s_n^\sharp$ and $t_n^\sharp$
going from $g_{n , *}\mathcal{O}_{Y_n}$ to $h_{n, *}\mathcal{O}_{R_n}$.
Here $g_n : Y_n = X' \amalg Z_n \to X$ and $h_n : R_n = Y_n \times_X Y_n \to X$
are the given morphisms. Let $\mathcal{I} \subset \mathcal{O}_X$ be the
coherent sheaf of ideals corresponding to $Z$. Then
$$
g_{n, *}\mathcal{O}_{Y_n} =
f_*\mathcal{O}_{X'} \times \mathcal{O}_X/\mathcal{I}^n
$$
Similarly, we have a decomposition
$$
R_n = X' \times_X X' \amalg X' \times_X Z_n \amalg
Z_n \times_X X' \amalg Z_n \times_X Z_n
$$
As $Z_n \to X$ is a monomorphism, we see that
$X' \times_X Z_n = Z_n \times_X X'$ and that this identification
is compatible with the two morphisms to $X$, with the two morphisms to
$X'$, and with the two morphisms to $Z_n$.
Denote $f_n : X' \times_X Z_n \to X$ the morphism to $X$.
Denote
$$
\mathcal{A} = \text{Equalizer}(
\xymatrix{
f_*\mathcal{O}_{X'} \ar@<1ex>[r] \ar@<-1ex>[r] &
(f \times f)_*\mathcal{O}_{X' \times_X X'}
}
)
$$
By the remarks above we find that
$$
\mathcal{A}_n =
\text{Equalizer}(
\xymatrix{
\mathcal{A} \times \mathcal{O}_X/\mathcal{I}^n \ar@<1ex>[r] \ar@<-1ex>[r] &
f_{n, *}\mathcal{O}_{X' \times_X Z_n}
}
)
$$
We have canonical maps
$$
\mathcal{O}_X \to \ldots \to \mathcal{A}_3 \to \mathcal{A}_2 \to \mathcal{A}_1
$$
of coherent $\mathcal{O}_X$-algebras. The statement of the lemma means that
for $n$ large enough there exists an $m \geq 0$ such that the image of
$\mathcal{A}_{n + m} \to \mathcal{A}_n$ is isomorphic to $\mathcal{O}_X$.
This we may check \'etale locally on $X$. Hence by Properties of Spaces,
Lemma \ref{spaces-properties-lemma-quasi-compact-affine-cover}
we may assume $X$ is an affine Noetherian scheme.

\medskip\noindent
Since $X_n \to X$ is an isomorphism over $U$ we see that the kernel
of $\mathcal{O}_X \to \mathcal{A}_n$ is supported on $|Z|$.
Since $X$ is Noetherian, the sequence of kernels
$\mathcal{J}_n = \Ker(\mathcal{O}_X \to \mathcal{A}_n)$ stabilizes
(Cohomology of Spaces, Lemma \ref{spaces-cohomology-lemma-acc-coherent}).
Say $\mathcal{J}_{n_0} = \mathcal{J}_{n_0 + 1} = \ldots = \mathcal{J}$.
By Cohomology of Spaces, Lemma
\ref{spaces-cohomology-lemma-power-ideal-kills-sheaf}
we find that $\mathcal{I}^t \mathcal{J} = 0$ for some $t \geq 0$.
On the other hand, there is an $\mathcal{O}_X$-algebra map
$\mathcal{A}_n \to \mathcal{O}_X/\mathcal{I}^n$
and hence $\mathcal{J} \subset \mathcal{I}^n$ for all $n$.
By Artin-Rees (Cohomology of Spaces, Lemma
\ref{spaces-cohomology-lemma-Artin-Rees}) we find that
$\mathcal{J} \cap \mathcal{I}^n \subset \mathcal{I}^{n - c}\mathcal{J}$
for some $c  \geq 0$ and all $n \gg 0$. We conclude that $\mathcal{J} = 0$.

\medskip\noindent
Pick $n \geq n_0$ as in the previous paragraph. Then
$\mathcal{O}_X \to \mathcal{A}_n$ is injective. Hence it now
suffices to find $m \geq 0$ such that the image of
$\mathcal{A}_{n + m} \to \mathcal{A}_n$ is equal
to the image of $\mathcal{O}_X$. Observe that $\mathcal{A}_n$
sits in a short exact sequence
$$
0 \to \Ker(\mathcal{A} \to f_{n, *}\mathcal{O}_{X' \times_X Z_n})
\to \mathcal{A}_n \to \mathcal{O}_X/\mathcal{I}^n \to 0
$$
and similarly for $\mathcal{A}_{n + m}$. Hence it suffices to show
$$
\Ker(\mathcal{A} \to f_{n + m, *}\mathcal{O}_{X' \times_X Z_{n + m}})
\subset
\Im(\mathcal{I}^n \to \mathcal{A})
$$
for some $m \geq 0$. To do this we may work \'etale locally on
$X$ and since $X$ is Noetherian we may assume that $X$ is
a Noetherian affine scheme. Say $X = \Spec(R)$ and $\mathcal{I}$
corresponds to the ideal $I \subset R$. Let $\mathcal{A} = \widetilde{A}$
for a finite $R$-algebra $A$. Let $f_*\mathcal{O}_{X'} = \widetilde{B}$
for a finite $R$-algebra $B$. Then $R \to A \subset B$ and these maps
become isomorphisms on inverting any element of $I$.

\medskip\noindent
Note that $f_{n, *}\mathcal{O}_{X' \times_X Z_n}$
is equal to $f_*(\mathcal{O}_{X'}/I^n\mathcal{O}_{X'})$
in the notation used in Cohomology of Spaces, Section
\ref{spaces-cohomology-section-theorem-formal-functions}.
By Cohomology of Spaces, Lemma
\ref{spaces-cohomology-lemma-ML-cohomology-powers-ideal}
we see that there exists a $c \geq 0$ such that
$$
\Ker(B \to \Gamma(X, f_*(\mathcal{O}_{X'}/I^{n + m + c}\mathcal{O}_{X'}))
$$
is contained in $I^{n + m}B$. On the other hand, as $R \to B$ is
finite and an isomorphism after inverting any element of $I$
we see that $I^{n + m}B \subset \Im(I^n \to B)$ for $m$ large enough
(can be chosen independent of $n$). This finishes the proof as $A \subset B$.
\end{proof}

\end{document}
